% Generated by scripts/assemble.py; corrections belong in source/corrections.json.
\examhead{197}{湘潭大学 2025 年数学分析试卷}{186}
% p187-b08; source PDF page 187
\blockanchor{p187-b08}
注. 部分题目顺序可能有误.\par
% p187-b09; source PDF page 187
\par\Needspace{5\baselineskip}
\blockanchor{p187-b09}
\majorquestion{一}求极限. 每题 10 分, 共 20 分.\par
% p187-b10; source PDF page 187
\blockanchor{p187-b10}
\question{1}$\lim_{n\to\infty}\left(1-\frac{1}{2^{2}}\right)\left(1-\frac{1}{3^{2}}\right)\cdots\left(1-\frac{1}{n^{2}}\right).$\par
% p187-b11; source PDF page 187
\blockanchor{p187-b11}
\question{2}$\lim_{x\to \infty}x^{\frac{3}{2}}\left(\sqrt{x + 1} +\sqrt{x - 1} -2\sqrt{x}\right).$\par
% p187-b12; source PDF page 187
\par\Needspace{5\baselineskip}
\blockanchor{p187-b12}
\majorquestion{二}求积分. 每题 10 分, 共 20 分.\par
% p187-b13; source PDF page 187
\blockanchor{p187-b13}
\question{1}$\int_0^{100\pi}\sqrt{1 - \cos 2x}\mathrm{d}x.$\par
% p187-b14; source PDF page 187
\blockanchor{p187-b14}
\question{2}$\lim_{n\to \infty}\int_0^1\frac{x^n}{1 + x}\mathrm{d}x.$\par
% p187-b15; source PDF page 187
\par\Needspace{5\baselineskip}
\blockanchor{p187-b15}
\majorquestion{三}(15 分) 设 $\{y_{n}\}$ 是严格单调增加的正无穷大量, 且 $\lim_{n\to \infty}\frac{x_n - x_{n - 1}}{y_n - y_{n - 1}} = 0$, 则 $\lim_{n\to \infty}\frac{x_n}{y_n} = 0$.\par
% p187-b16; source PDF page 187
\par\Needspace{5\baselineskip}
\blockanchor{p187-b16}
\majorquestion{四}(10 分) 证明不等式: $x - \frac{x^2}{2} < \ln(1 + x) < x, x > 0$.\par
% p187-b17; source PDF page 187
\par\Needspace{5\baselineskip}
\blockanchor{p187-b17}
\majorquestion{五}（15分）讨论反常积分的敛散性： $\int_0^{+\infty}\frac{\ln(1 + x)}{x^p}\mathrm{d}x,p\in \mathbb{R}.$\par
% p187-b18; source PDF page 187
\par\Needspace{5\baselineskip}
\blockanchor{p187-b18}
\majorquestion{六}(10 分) 计算重积分 $\iint_{D} xy^{2} \, \mathrm{d}x \, \mathrm{d}y$, 其中 $D$ 为抛物线 $y^{2} = 2px$ 和直线 $x = \frac{p}{2} (p > 0)$ 所围的区域.\par
% p187-b19; source PDF page 187
\par\Needspace{5\baselineskip}
\blockanchor{p187-b19}
\majorquestion{七}(15 分) 设 $f(x)$ 在区间 $I$ 上为下凸函数, 证明: $f(x)$ 在区间 $I$ 的任一闭子区间上有界.\par
% p187-b20; source PDF page 187
\par\Needspace{5\baselineskip}
\blockanchor{p187-b20}
\majorquestion{八}(15 分) 设函数 $\varphi(x)$ 在 $[0, +\infty)$ 上存在二阶连续的导数, 若 $\lim_{x \to +\infty} \varphi(x)$ 存在, 且 $\varphi''(x)$ 有界. 试证: $\lim_{x \to +\infty} \varphi'(x) = 0$.\par
% p187-b21; source PDF page 187
\par\Needspace{5\baselineskip}
\blockanchor{p187-b21}
\majorquestion{九}(15 分) 设 $f(x) = \sum_{n=0}^{\infty} \frac{1}{2^n + x}$.\par
% p187-b22; source PDF page 187
\blockanchor{p187-b22}
(1) (9 分) 证明: $f(x)$ 在 $[0, +\infty)$ 上可导, 且一致连续.\par
% p187-b23; source PDF page 187
\blockanchor{p187-b23}
(2) (6 分) 证明: 反常积分 $\int_{0}^{+\infty} f(x) \, \mathrm{d}x$ 发散.\par
% p187-b24; source PDF page 187
\par\Needspace{5\baselineskip}
\par\Needspace{6\baselineskip}
\blockanchor{p187-b24}
\majorquestion{十}(15 分) 计算曲面积分\par
\nopagebreak[4]
% p187-b25; source PDF page 187
\blockanchor{p187-b25}
\displayformula{\iint_ {S} \left(y - x ^ {3}\right) \mathrm{d} y \mathrm{d} z + \left(y ^ {3} - z\right) \mathrm{d} z \mathrm{d} x + 2 \mathrm{d} x \mathrm{d} y.}
% p187-b26; source PDF page 187
\blockanchor{p187-b26}
其中 $S$ 是锥面 $z^2 = x^2 + y^2 (0 \leq z \leq 2)$, 方向取下侧.\par
